problem:  
Let \(G\) be a connected simply connected **2-step nilpotent Lie group** with Lie algebra \(\mathfrak{g}\).  
Every left-invariant Lorentzian metric \(g\) on \(G\) corresponds to a symmetric bilinear form \(B\) of signature \((1, n-1)\) on \(\mathfrak{g}\). The geodesic equations reduce to the **Euler equation** on the Lie algebra:  
\[
\dot{u} = -\,\operatorname{ad}(u)^{\top}u, \quad u(t)\in \mathfrak{g},
\]
with \(\operatorname{ad}(u)^{\top}\) defined using \(B\).  

It is known that for Riemannian left-invariant metrics on 2-step nilpotent groups, geodesic completeness always holds. By contrast, in the Lorentzian case completeness can fail in subtle ways.  

**Problem:**  
Find precise **algebraic conditions** on a pair \((\mathfrak{g}, B)\), where \(\mathfrak{g}\) is 2-step nilpotent and \(B\) a Lorentzian inner product, that are **necessary and sufficient** for the associated left-invariant metric on \(G\) to be geodesically complete.  

Equivalently: characterize those Lorentzian forms \(B\) on a fixed 2-step nilpotent Lie algebra \(\mathfrak{g}\) such that all trajectories of the Euler equation above are complete for all initial data \(u(0)\in \mathfrak{g}\).  

A finer version of the problem asks:  
Is geodesic completeness a **stable property** under small deformations of \(B\) among Lorentzian forms, when \(\mathfrak{g}\) is 2-step nilpotent?

Why is it a "good" problem:  
This problem strikes directly at the interface of **Lorentzian geometry**, **Lie theory**, and **dynamical systems**. The reduction of the geodesic equation to an algebraic quadratic vector field on \(\mathfrak{g}\) transforms a geometric completeness problem into a deep **global dynamics question** on an algebraic ODE system—one sensitive to the indefinite signature. It naturally generalizes the known Riemannian completeness of nilpotent groups to the much more intricate Lorentzian setting, probing whether algebraic nilpotency still enforces dynamical completeness or whether the indefinite metric breaks that stability. The question’s resolution would require synthesizing tools from geometric analysis, the qualitative theory of ODEs, and the structure theory of Lie algebras, leading possibly to a new **algebraic classification of completeness phenomena**. The problem is simple to state—“When is a Lorentzian left-invariant metric on a 2-step nilpotent group geodesically complete?”—yet cuts to a foundational issue linking curvature, algebraic structure, and global dynamics.